\documentclass[twocolumn]{aastex631}
\usepackage{amsmath}
\usepackage{booktabs}
\shorttitle{SDSS target vector for feedback-model validation}
\shortauthors{NebulaMind Research Autopilot}
\begin{document}

\title{SDSS target vector for feedback-model validation}
\author{NebulaMind Research Autopilot}
\affiliation{NebulaMind Astrophysics Collaboration, San Francisco, CA 94107, USA}
\correspondingauthor{NebulaMind Research Autopilot}
\email{autopilot@nebulamind.ai}

\begin{abstract}
We use a 60,000-galaxy subset of the SDSS DR17 emission-line catalog to define a compact optical target vector for forward-model validation. The pilot records quenched fraction, optical AGN incidence, and color versus mass/redshift across 15 mass-redshift cells with $n \geq 50$; across mass bins, quenched fractions span 0.005--0.729 and optical AGN fractions span 0.003--0.520. It remains an empirical denominator study rather than a direct simulation comparison.
\end{abstract}

\keywords{galaxies: evolution --- galaxies: active --- galaxies: star formation --- surveys --- methods: data analysis}
\software{Astropy, SciPy, NumPy, Matplotlib, pandas}

\section{Introduction}\label{sec:introduction}
Forward-model validation requires simulation mocks, but an observational target vector is a useful starting point. Here we present the SDSS DR17 emission-line sample as a compact optical baseline and restrict the manuscript to directly measured quantities. Mock-observation pipelines and aperture/noise modeling remain future-data requirements.


\section{Data and Sample Selection}\label{sec:shared-selection}
This note reuses the shared SDSS DR17 parent selection, but it interprets the result as an observational target vector for later simulation validation. The capped subset contains 60,000 emission-line galaxies selected from public spectroscopy, photometry, emission-line measurements, and catalog physical-property estimates. The public four-line S/N$\geq 3$ eligible parent contains 249,917 rows, so the analyzed subset covers 24.0\% of that parent. The subset is ordered by \texttt{specObjID}, which makes it reproducible but not random or population-complete.

\begin{deluxetable*}{lrrr}
\tabletypesize{\scriptsize}
\tablecaption{Shared SDSS DR17 selection cascade used before paper-specific quantities.\label{tab:selection-cascade}}
\tablehead{\colhead{Selection stage} & \colhead{Public DR17 rows} & \colhead{Cached rows} & \colhead{Retention vs. spectro-z parent (fraction)}}
\startdata
SpecObj GALAXY, 0.02<z<0.12 & 501,060 & -- & 1.000 \\
plus galSpecInfo/PhotoObj/galSpecExtra and mass/sSFR bounds & 416,554 & -- & 0.831 \\
plus galSpecLine join & 416,554 & -- & 0.831 \\
four BPT lines positive with positive errors & 373,445 & 60,000 & 0.745 \\
four BPT lines S/N$\geq 3$ & 249,917 & 60,000 & 0.499 \\
four BPT lines S/N$\geq 5$ & 176,523 & 42,446 & 0.352 \\
four BPT lines S/N$\geq 10$ & 91,768 & 22,311 & 0.183 \\
\enddata
\tablecomments{Counts are read-only public SDSS DR17 count queries plus the cached local CSV. Cached rows are shown only where the cache applies.}
\end{deluxetable*}

The four-line requirement is strongly selection dependent. In the public counts, S/N$\geq 3$ keeps 33.6\% of the $-12<\log {\rm sSFR}<-11$ parent bin but 94.9\% of the $-10<\log {\rm sSFR}<-9.5$ bin. Therefore every incidence, quenching, density, gas-denominator, or target-vector statement in this analysis is conditional on the four-line emission-line selection.

Local subset versus public catalog marginal checks found no redshift, stellar-mass, or sSFR bin with a subset-minus-public fraction difference above 5 percentage points. The largest absolute differences were 2.03 percentage points in redshift, -1.63 percentage points in stellar mass, and -0.58 percentage points in sSFR. This is a representativeness diagnostic only; it does not make the capped cache random or complete.


\section{Measurements}\label{sec:measurements}
The row-level measurements include redshift, stellar mass, catalog specific star-formation rate, model colors, and four optical line fluxes/errors. BPT labels and all derived denominators are recomputed locally from the cached table \citep{baldwin1981,kewley2001,kauffmann2003bpt,kewley2006}. Catalog SFR/sSFR values are treated as low-redshift SDSS physical-property estimates rather than direct resolved gas or feedback measurements \citep{sdssdr17,brinchmann2004,york2000}.
Unless otherwise noted, quoted fraction uncertainties are binomial counting uncertainties from the stated sample sizes, and bracketed intervals are bootstrap confidence intervals.


\section{Optical target vector for simulation validation}\label{sec:topic-result}
We define a compact SDSS target vector of quenched fraction, optical AGN incidence, and color versus mass/redshift for forward-model validation. The result is an observed optical baseline rather than a full physical-feedback test.

We define 15 mass-redshift cells with $n \geq 50$ as a compact validation vector; the cell grid spans mass bins 8.0--9.5, 9.5--10.0, 10.0--10.5, 10.5--11.0, and 11.0--12.5 crossed with redshift bins 0.02--0.05, 0.05--0.08, and 0.08--0.12. Across mass bins, quenched fractions span 0.005--0.729 and optical AGN fractions span 0.003--0.520. The observed target vector is useful for simulation forward modelling, but it still requires mock-observation pipelines before any model comparison can be claimed.


\begin{figure}
\centering
\includegraphics[width=\columnwidth]{../figures/fig-topic.pdf}
\caption{SDSS DR17 optical denominator/proxy diagnostic for the feedback-model validation target vector. The figure maps quenched fractions and optical AGN incidence across 15 mass-redshift cells for simulation forward-modeling, spanning mass bins 8.0--9.5 through 11.0--12.5 and redshift bins 0.02--0.05 through 0.08--0.12.}
\label{fig:topic}
\end{figure}

\section{Interpretation and missing observables}\label{sec:missing}
This SDSS-only pilot is intentionally limited to optical quantities. A full proposal requires simulation mocks passed through the SDSS/MaNGA/ALMA/X-ray/radio selection functions and aperture/noise models.

Simulation suites such as TNG, EAGLE, and SIMBA define the future comparison problem; the iMaNGA observational catalog provides a complementary benchmark for mock-observation work. No simulation mock has been forward-modelled or ranked in this pilot \citep{tng2019,eagle2015,simba2019,imanga2023,donnari2021,dubois2013,dubois2016}.


\section{Data Availability}\label{sec:data-avail}
The SDSS DR17 spectroscopy, photometry, emission-line measurements, and catalog quantities used here are publicly available through the SDSS data releases and associated catalog papers cited below. A local subset and manifest are retained in the project repository for reproducibility and are available from the corresponding author upon reasonable request.

\section{Conclusion}\label{sec:conclusion}
We define 15 mass-redshift cells with $n \geq 50$ as a compact validation vector, spanning mass bins 8.0--9.5, 9.5--10.0, 10.0--10.5, 10.5--11.0, and 11.0--12.5 across redshift bins 0.02--0.05, 0.05--0.08, and 0.08--0.12. Quenched fractions span 0.005--0.729 and optical AGN fractions span 0.003--0.520. This observed target vector is useful for simulation forward modelling, but it still requires mock-observation pipelines before any model comparison can be claimed.

\acknowledgments
We thank the SDSS collaboration. This manuscript uses public SDSS DR17 data only.


\begin{thebibliography}{}
\bibitem[Abdurro'uf et al.(2022)]{sdssdr17} Abdurro'uf, Accetta, K., Aerts, C., et al. 2022, ApJS, 259, 35
\bibitem[Baldwin et al.(1981)]{baldwin1981} Baldwin, J.~A., Phillips, M.~M., \& Terlevich, R. 1981, PASP, 93, 5
\bibitem[Brinchmann et al.(2004)]{brinchmann2004} Brinchmann, J., Charlot, S., White, S.~D.~M., et al. 2004, MNRAS, 351, 1151
\bibitem[Kauffmann et al.(2003a)]{kauffmann2003bpt} Kauffmann, G., Heckman, T.~M., Tremonti, C., et al. 2003a, MNRAS, 346, 1055
\bibitem[Kewley et al.(2001)]{kewley2001} Kewley, L.~J., Dopita, M.~A., Sutherland, R.~S., Heisler, C.~A., \& Trevena, J. 2001, ApJ, 556, 121
\bibitem[Kewley et al.(2006)]{kewley2006} Kewley, L.~J., Groves, B., Kauffmann, G., \& Heckman, T. 2006, MNRAS, 372, 961
\bibitem[York et al.(2000)]{york2000} York, D.~G., Adelman, J., Anderson, J.~E., Jr., et al. 2000, AJ, 120, 1579
\bibitem[Dav\'e et al.(2019)]{simba2019} Dav\'e, R., Angl\'es-Alc\'azar, D., Narayanan, D., et al. 2019, MNRAS, 486, 2827
\bibitem[Donnari et al.(2021)]{donnari2021} Donnari, M., Pillepich, A., Nelson, D., et al. 2021, MNRAS, 506, 4760
\bibitem[Dubois et al.(2013)]{dubois2013} Dubois, Y., Gavazzi, R., Peirani, S., \& Silk, J. 2013, MNRAS, 433, 3297
\bibitem[Dubois et al.(2016)]{dubois2016} Dubois, Y., Peirani, S., Pichon, C., et al. 2016, MNRAS, 463, 3948
\bibitem[Nanni et al.(2023)]{imanga2023} Nanni, L., Thomas, D., Trayford, J., et al. 2023, MNRAS, 518, 2605
\bibitem[Nelson et al.(2019)]{tng2019} Nelson, D., Springel, V., Pillepich, A., et al. 2019, Computational Astrophysics and Cosmology, 6, 2
\bibitem[Schaye et al.(2015)]{eagle2015} Schaye, J., Crain, R.~A., Bower, R.~G., et al. 2015, MNRAS, 446, 521
\end{thebibliography}

\end{document}
